\section{Affine improvement under subdivision}
\label{sec:subdivision}

The subharmonic profiles of Section~\ref{sec:profiles} admit a useful
local affine correction.  Averaged over a sufficiently fine square
grid, the remaining growth decreases faster than the side length.
We first establish this fact for profiles and then transfer it to
eigenfunctions.

\begin{samepage}
\begin{definition}[Growth and affine excess]
\label{def:growth}
Fix an isothermal coordinate region in which
\(\Delta u+k^2p u=0\), and work in a relatively compact subset on which
\(p\) and \(\nabla p\) have fixed uniform bounds.  Set \(R=1000\).
Choose a root square \(Q_0\), with side \(s_0\le1\) and center
\(x_{Q_0}\), such that
\(\overline{\D_{2Rs_0}(x_{Q_0})}\) lies in this subset.
For an integer \(A\ge2\), subdivide each square into \(A^2\) congruent
squares and write \(q\prec Q\) when \(q\) is an immediate child of
\(Q\).  For a square \(Q\) of side \(s=s(Q)\) and center \(x_Q\), put
\(K_Q=ks\).  For a real vector \(b\in\R^2\), define
\begin{equation}
\begin{aligned}
 E(Q,b)
 &=\log
 \frac{\norm{e^{-kb\cdot(x-x_Q)}u}_{L^2(\D_{Rs}(x_Q))}}
      {\norm{e^{-kb\cdot(x-x_Q)}u}_{L^2(\D_s(x_Q))}},\\
 m(Q,b)&=\frac{E(Q,b)}{K_Q},
 \qquad N(Q)=m(Q,0).
\end{aligned}
\label{eq:growth}
\end{equation}
All planar norms here and below are unnormalized Euclidean norms.
An average over children assigns weight \(A^{-2}\) to each child.
\end{definition}
\end{samepage}

A square \(Q\) is called \emph{admissible} if \(s(Q)\le1\) and
\(\overline{\D_{Rs(Q)}(x_Q)}\) lies in the fixed coordinate region
where the coefficient bounds hold.  All descendants of the chosen root
are admissible.  The stronger \(2R\) margin at the root reserves room
for nearby squares outside its subdivision tree.

Proposition~\ref{prop:initial-growth} ensures that the local norms in
\eqref{eq:growth} are nonzero.  In particular, \(E\), \(m\), and \(N\)
are finite and nonnegative, since the disks are nested.  Comparing the
exponential weight on the two disks gives
\begin{equation}
 N(Q)\le m(Q,b)+(R+1)\abs{b}.
\label{eq:untilted-comparison}
\end{equation}
Indeed the outer untilted norm is at most \(e^{KR\abs b}\) times its
tilted norm, whereas the inner untilted norm is at least
\(e^{-K\abs b}\) times its tilted norm, where \(K=K_Q\).
The same definitions and estimates apply to finite subtrees of an
extended translated grid consisting of admissible squares.

The normalization by \(K_Q\) explains why subdivision alone is
insufficient.  Put \(\rho=A^{-1}\).  An affine profile with slope
\(c\ne0\) has a difference of suprema equal to
\((R-1)\rho|c|\) between concentric disks of radii \(R\rho\) and
\(\rho\).  Division by the child frequency \(K_q=\rho K_Q\)
cancels this first-order gain.  We will subtract the local affine
part to obtain an average profile excess of order
\(\rho^2(1+|\log\rho|)\), while keeping the average correction slope
bounded.  The latter bound controls the cost of changing weights
through successive subdivisions.

\begin{lemma}[Uniform potential decomposition]
\label{lem:potential-decomposition}
Let \(V\) be subharmonic on \(\D_R\), with
\[
 V\le0,
 \qquad \sup_{\overline{\D_1}}V\ge-1.
\]
Then \(\norm{V}_{L^1(\D_8)}\) and the mass of its Riesz measure
\((2\pi)^{-1}\Delta V\) on \(\D_5\) are uniformly bounded.
There are a nonnegative measure \(\nu\) on \(\D_4\), of uniformly
bounded mass, and a harmonic function \(h\) on \(\D_3\) such that
\[
 V(x)=h(x)+\int_{\D_4}\log\abs{x-z}\,d\nu(z),
 \qquad x\in\D_3.
\]
The first and second derivatives of \(h\) on \(\D_2\) are uniformly
bounded.  All bounds are independent of \(V\), and the representation
holds for the pointwise subharmonic representatives, including values
equal to \(-\infty\).
\end{lemma}

\begin{proof}
Upper semicontinuity supplies a point
\(x_*\in\overline{\D_1}\) with \(V(x_*)\ge-1\).  The disk
\(\D_{10}(x_*)\) contains \(\D_8\) and is contained in \(\D_R\).
In particular, \(V\) is nontrivial and locally integrable.
Since \(V\le0\), the submean inequality gives
\[
 \int_{\D_8}\abs{V}
 \le -\int_{\D_{10}(x_*)}V
 \le100\pi.
\]
Its distributional Laplacian is a nonnegative Radon measure.
Choose a fixed nonnegative
smooth cutoff \(\psi\), supported in \(\D_8\) and equal to one on
\(\D_5\).  If \(\eta=(2\pi)^{-1}\Delta V\), then
\[
 \eta(\D_5)
 \le\int\psi\,d\eta
 =\frac1{2\pi}\int V\Delta\psi
 \le C\norm{V}_{L^1(\D_8)}.
\]
Let \(\nu\) be the restriction of \(\eta\) to \(\D_4\), and set
\(P(x)=\int_{\D_4}\log\abs{x-z}\,d\nu(z)\).
The elementary bound
\[
 \sup_{z\in\D_4}\int_{\D_3}\bigl|\log\abs{x-z}\bigr|\,dx<\infty
\]
shows that \(\norm{P}_{L^1(\D_3)}\le C\nu(\D_4)\).
The logarithmic fundamental solution identity
\(\Delta\log\abs{x-z}=2\pi\delta_z\) implies
\(\Delta(V-P)=0\) on \(\D_3\).  Thus \(V-P\) agrees almost everywhere
with a harmonic function \(h\).  Interior harmonic estimates, applied
to its bounded \(L^1(\D_3)\) norm, bound \(\nabla h\) and \(\nabla^2h\)
on \(\D_2\).  Finally, the logarithmic potential of a nonnegative
measure has its usual subharmonic representative.  Both \(V\) and
\(h+P\) are recovered pointwise as limits of their disk averages.
Their almost-everywhere equality therefore gives the asserted
pointwise representation.
\end{proof}

\begin{lemma}[Affine correction of a profile]
\label{lem:affine-profile}
Let \(V\), \(h\), and \(\nu\) be as in
Lemma~\ref{lem:potential-decomposition}.  Subdivide the unit square
\([-1/2,1/2]^2\) into \(A^2\) squares, where \(A>10R\), put
\(\rho=A^{-1}\), and denote their centers by \(y\).  Define
\[
 c_y=\nabla h(y)
   +\int_{\substack{z\in\D_4\\\abs{z-y}>2R\rho}}
        \frac{y-z}{\abs{y-z}^2}\,d\nu(z).
\]
Then
\begin{equation}
\begin{aligned}
 &\sup_{\overline{\D_{R\rho}(y)}}(V(x)-c_y\cdot x)
     -\sup_{\D_\rho(y)}(V(x)-c_y\cdot x)\\
 &\hspace{1cm}\le
 C\left(\rho^2+\int_{\D_4}
      \min\left\{1,\frac{\rho^2}{\abs{z-y}^2}\right\}\,d\nu(z)
   \right).
\end{aligned}
\label{eq:affine-error}
\end{equation}
At \(z=y\) the minimum in this formula is interpreted as one.
Moreover,
\[
 \av_y\abs{c_y}\le C,
 \qquad
 \av_y\left[\text{right-hand side of \eqref{eq:affine-error}}\right]
 \le C\rho^2(1+\abs{\log\rho}).
\]
The constants may depend on \(R\), but are independent of \(A\) and
\(V\).
\end{lemma}

\begin{proof}
All disks \(\overline{\D_{R\rho}(y)}\) lie in \(\D_2\).  The vector
\(c_y\) is finite because its kernel is truncated away from its
singularity.  Put \(W_y(x)=V(x)-c_y\cdot x\), and lower bound its
supremum on the inner disk by its disk average there.  For the
harmonic term, this average is its value at \(y\), after the affine
part is accounted for.  Taylor's theorem bounds the corresponding
outer deviation by \(C\rho^2\).

For a far source \(\abs{z-y}>2R\rho\), the logarithm is harmonic
on the outer disk and its Hessian is bounded there by
\(C\abs{z-y}^{-2}\).  Its inner disk average is
\(\log\abs{y-z}\), and the retained gradient in \(c_y\) removes
its first-order term.  Its contribution to the outer deviation is
therefore at most \(C\rho^2/\abs{z-y}^2\).

For a near source \(\abs{z-y}\le2R\rho\), rescale by \(\rho\).
Writing \(t=(z-y)/\rho\), the disk-average formula
\[
 \frac1\pi\int_{\D_1}\log\abs{\xi-t}\,d\xi
 =
 \begin{cases}
  (\abs{t}^2-1)/2,&\abs{t}\le1,\\
  \log\abs{t},&\abs{t}\ge1
 \end{cases}
\]
has lower bound \(-1/2\).  On the larger disk the corresponding
logarithm is at most \(\log(3R)\), since \(\abs t\le2R\).
The common term \(\log\rho\) cancels between these two quantities.
Thus the difference between the outer upper bound and the inner
average is at most \(\log(3R)+1/2\), uniformly in the source position.
This calculation permits an atom at \(y\), or anywhere else in
the inner disk: the logarithm remains integrable in the averaging
variable.  At an outer point where a potential is \(-\infty\), an
upper bound causes no difficulty.  Because the measure is
nonnegative, the kernel bounds may be integrated.  On the near
region,
\[
 \min\left\{1,\frac{\rho^2}{\abs{z-y}^2}\right\}
 \ge\frac1{4R^2}.
\]
Combining the near and far bounds proves
\eqref{eq:affine-error}.

For completeness, fix any \(z\in\D_4\).  The number of grid centers
within distance \(r\) of \(z\) is at most
\(C(1+r^2/\rho^2)\).  There are \(\rho^{-2}\) centers in total.
The centers at distance at most \(\rho\) contribute \(C\rho^2\)
to the averaged minimum kernel.  In each dyadic annulus
\(2^\ell\rho<\abs{z-y}\le2^{\ell+1}\rho\), the number of centers
is at most \(C2^{2\ell}\), whereas the kernel is at most
\(2^{-2\ell}\).  All distances are bounded by five, so only
\(C(1+\abs{\log\rho})\) annuli occur.  Hence, uniformly in \(z\),
\[
 \av_y\min\left\{1,\frac{\rho^2}{\abs{z-y}^2}\right\}
 \le C\rho^2(1+\abs{\log\rho}).
\]
The same annuli, for the truncated gradient kernel, give
\[
 \rho^2\sum_{\substack{y\\\abs{z-y}>2R\rho}}
 \frac1{\abs{z-y}}
 \le C.
\]
Indeed an annulus of radius comparable to \(2^\ell\rho\)
contributes at most \(C2^\ell\rho\), and these radii sum to a
fixed constant.  Integrating the two bounds against the bounded
measure \(\nu\), and using the bound for \(\nabla h\), proves
both averaged estimates.
\end{proof}

\begin{proposition}[One subdivision step]
\label{prop:subdivision}
There are an integer \(A>10R\) and constants \(C_0,C_1>0\), depending
only on the fixed coordinate coefficient bounds and \(R\), with the
following property.  Let \(Q\) be an admissible square of side
\(s\le1\), let \(K=ks>0\), and let \(b\in\R^2\) and \(S>0\).
Put \(B=Kb\).  If
\begin{equation}
 S\ge C_0,\qquad E(Q,b)\le S,\qquad
 \frac{K}{S+\abs B}\le\frac1{C_0},\qquad
 \frac{K^2s}{S(S+\abs B)}\le\frac1{C_0},
\label{eq:subdivision-hypotheses}
\end{equation}
then real vectors \(b_q\) can be chosen for the children \(q\prec Q\)
so that
\begin{equation}
 \av_{q\prec Q}\abs{b_q-b}\le C_1\frac SK,
 \qquad
 \av_{q\prec Q}m(q,b_q)\le\frac12\frac SK.
\label{eq:subdivision-estimate}
\end{equation}
The constants are independent of the eigenfunction, the scale,
the depth of a tree, and any lower bound for \(K\).
\end{proposition}

\begin{proof}
First choose \(A>10R\), using only the uniform estimates of
Lemma~\ref{lem:affine-profile}, so large that, with \(\rho=A^{-1}\),
\[
 C\rho^2(1+\abs{\log\rho})<\frac\rho4.
\]
Choose \(C_1\) strictly larger than the uniform bound for
\(\av_y\abs{c_y}\) in that lemma.  We prove that some finite
\(C_0\) works for these fixed choices.

Suppose otherwise.  There is a sequence of failures of the
conclusion for which the hypotheses hold with \(C_0=j\).
Use subscripts \(j\) for the associated quantities and rescale each
parent square to the unit square centered at zero:
\[
 U_j(X)=u_j(x_{Q_j}+s_jX),
 \qquad p_j(X)=p(x_{Q_j}+s_jX).
\]
The rescaled equation is
\(\Delta U_j+K_j^2p_jU_j=0\), and
\(\norm{p_j}_\infty\le C\),
\(\norm{\nabla p_j}_\infty\le Cs_j\).
With \(B_j=K_jb_j\) and \(D_j=S_j+\abs{B_j}\), the failure
hypotheses imply
\[
 S_j\longrightarrow\infty,\qquad
 \frac{K_j}{D_j}\longrightarrow0,\qquad
 \frac{K_j^2\norm{\nabla p_j}_\infty}{S_jD_j}
 \longrightarrow0.
\]
No restriction on \(\abs{B_j}/S_j\) has been imposed.

Normalize the tilted measures on \(\D_R\) by setting
\[
 d\mu_j(X)=
 \frac{e^{-2B_j\cdot X}\abs{U_j(X)}^2\,dX}
 {\displaystyle\int_{\D_R}e^{-2B_j\cdot X}\abs{U_j(X)}^2\,dX}.
\]
Lemma~\ref{lem:profile-extraction} and
Theorem~\ref{thm:subharmonic-profile} give, after passage to a
subsequence, a subharmonic profile \(V\), possibly identically
\(-\infty\), on \(\D_R\).  It satisfies \(V\le0\), while
\[
 \mu_j(\D_1)=e^{-2E(Q_j,b_j)}\ge e^{-2S_j}.
\]
The compact upper bound in \eqref{eq:profile-bounds}, applied to
\(\overline{\D_1}\), therefore gives
\(\sup_{\overline{\D_1}}V\ge-1\).  In particular, \(V\) is
nontrivial and the two preceding lemmas apply.

For each of the fixed, finitely many child centers \(y\), take
the vector \(c_y\) furnished by Lemma~\ref{lem:affine-profile} for
this limiting profile.  On the corresponding child \(q_j\), set
\[
 b_{q_j}=b_j+\frac{S_j}{K_j}c_y.
\]
In parent coordinates the additional weight is
\(e^{-2S_jc_y\cdot X}\).  To see the cancellation of recentering
constants, write
\[
 -K_jb_{q_j}\cdot(X-y)
 =-B_j\cdot X-S_jc_y\cdot X
   +(B_j+S_jc_y)\cdot y.
\]
The last term, as well as the measure normalization and the
coordinate Jacobian, cancels in the ratio defining \(E(q_j,b_{q_j})\).
Consequently this excess is one half the logarithm of
\[
 \frac{\displaystyle\int_{\D_{R\rho}(y)}
                  e^{-2S_jc_y\cdot X}\,d\mu_j(X)}
      {\displaystyle\int_{\D_\rho(y)}
                  e^{-2S_jc_y\cdot X}\,d\mu_j(X)}.
\]
Apply the weighted compact upper bound to the numerator, using
\(\overline{\D_{R\rho}(y)}\), and the weighted open lower bound
to the denominator.  This yields
\[
 \limsup_j\frac{E(q_j,b_{q_j})}{S_j}
 \le
 \sup_{\overline{\D_{R\rho}(y)}}(V(X)-c_y\cdot X)
 -\sup_{\D_\rho(y)}(V(X)-c_y\cdot X).
\]
The inner supremum is finite: a nontrivial subharmonic function is
locally integrable, so it cannot equal \(-\infty\) throughout an
open disk, and it is bounded above on compact subsets.  Both
suprema in the displayed difference are therefore finite.

There are only \(A^2\) children, and \(A\) has already been fixed.
Averaging the last inequality and invoking
Lemma~\ref{lem:affine-profile} gives
\[
 \limsup_j\av_{q_j\prec Q_j}
       \frac{E(q_j,b_{q_j})}{S_j}<\frac\rho4.
\]
Since \(K_{q_j}=\rho K_j\), it follows, for all sufficiently large
\(j\), that
\[
 \av_{q_j\prec Q_j}m(q_j,b_{q_j})
 \le\frac12\frac{S_j}{K_j}.
\]
At the same time,
\[
 \av_{q_j\prec Q_j}\abs{b_{q_j}-b_j}
 =\frac{S_j}{K_j}\av_y\abs{c_y}
 \le C_1\frac{S_j}{K_j}.
\]
Thus these choices satisfy both required inequalities, contrary
to the defining failure of the sequence.  This proves the
proposition.  In particular, \(A\) and \(C_1\) have been fixed
before \(C_0\), and no positive lower bound for \(K\) was used.
\end{proof}
